\subsection{Full Factorial and Search-Retention Evidence}
\label{app:factorial-details}

Section~\ref{M-subsec:grounding-followup} distinguishes the historical 40-versus-7 all-key result from the controlled 49-query study. The latter's 277 scalar slots and reference-derived template inventory are fixed within each decoder comparison. The known excluded query shares a delexicalized text group with earlier training/development data; grouping uses transitive five-gram Jaccard similarity at least 0.80. Exclusion addresses that known crossing, not prior benchmark exposure or unrecognized semantic duplicates.

The design crosses two candidate constructors, three scorers, two decoders
and two mention-reuse policies (24 cells). Decoder pairs share the exact
same edge-score matrix, canonical lexical slot order, and zero-score null
option. Scorers are an explicit ordinal encoding of the original type
preference, the frozen contextual scorer, and a subtraction of only its
percent/currency marker bonuses. The ordinal control introduces a documented
cardinalization and does not reproduce every historical assignment.
One-to-one assignment is a diagnostic restriction rather than a physical
law. With unrestricted reuse and additive scores, independent per-slot
maximization is already globally optimal; every corresponding local/global
output agrees. Table~\ref{tab:grounding-factorial-diagnostic} shows the
one-to-one controls; all 24 cells, candidate identities, matrices and
per-query predictions are retained in the reproducibility artifacts.

\begin{table}[t]
\centering
\caption{One-to-one factorial controls on the same 49 queries and 277 gold
scalar slots. Entries show greedy $\rightarrow$ exact maximum-sum assignment.
E/T denote numerical exact/tolerance-correct counts. Full instances also
require compatible types and every gold scalar key. These controls are
distinct from the historical 331-query configurations.}
\label{tab:grounding-factorial-diagnostic}
\small
\setlength{\tabcolsep}{3pt}
\begin{tabular}{@{}llcccc@{}}
\toprule
Pool & Scorer & E slots & T slots & E full & T full\\
\midrule
Typed & Ordinal & $71\to71$ & $87\to86$ & $3\to3$ & $3\to3$\\
Typed & Context & $106\to136$ & $117\to145$ & $3\to4$ & $5\to5$\\
Typed & No marker bonus & $97\to137$ & $111\to146$ & $3\to4$ & $5\to5$\\
Rich & Ordinal & $67\to67$ & $81\to81$ & $3\to3$ & $3\to3$\\
Rich & Context & $100\to133$ & $111\to142$ & $3\to4$ & $5\to5$\\
Rich & No marker bonus & $97\to135$ & $111\to144$ & $3\to4$ & $5\to5$\\
\bottomrule
\end{tabular}
\end{table}

With contextual scores, exact assignment increases exact numerical slot
counts by 30/277 (typed pool) and 33/277 (rich pool). Existing paired
query-group bootstrap intervals ($B=2000$, seed 42) are respectively
$[3.14,17.95]$ and $[4.07,18.84]$ percentage points. These are marginal,
exploratory intervals. The full-tolerance count remains 5/49 for both
decoders, with one gained and one lost query; the paired interval is
$[-6.12,6.12]$ percentage points. A separate retrospective exact groupwise
method-swap analysis covers every one-factor contrast and four numerical
endpoints, using a single Holm family of 240 tests
\cite{holm1979followup}. Neither contextual-score decoder contrast survives
that familywise adjustment (adjusted $p=1$). Some contextual-versus-ordinal
slot comparisons survive; no whole-instance endpoint does. These results
support descriptive component differences, not a confirmed whole-instance
improvement or attribution of all 33 historical additional successes to
search. The recorded sub-millisecond cell timings include metric and
instrumentation overhead and are not clean algorithm-speed comparisons.

\medskip\noindent\textbf{Candidate availability, pruning and objective quality.}
The earlier 50-query joint-grounding pilot selects 3 tolerance-correct
instances with either local or joint scoring. In the saved common pools,
32/50 permit a criterion-correct complete mapping, but only 8/50 retain one
in the top ten. Reconciliation with identical saved beam outputs separates
14 cases losing correct prefixes, 6 generating then pruning a complete
correct mapping, 4 losing it at the final top-ten cut, 5 retaining it below
the selected mapping, and 3 selecting it. The other 18 cases lack a complete
compatible mapping in the pool or inventory. Width and ranking limitations
must therefore be distinguished from candidate absence. Optimality claims
about beam methods also depend on their search restriction
\cite{huang2017beamfollowup}; that neural-generation result is methodological
context, not a guarantee for this grounding algorithm.

An exact MILP diagnostic holds the cached local/pairwise objective and
domains fixed. Among the 32 pool-compatible cases, 23 have a strictly
higher-scoring incorrect optimum than their best criterion-correct mapping;
6 have tied correct and incorrect optima, and 3 select a correct optimum.
The correctness-restricted solves are oracle diagnostics, never inference
procedures. Exact mention-based top-ten enumeration retains a correct
mapping in 9/50 cases versus 8/50 for the beam, with tied-optimum membership
qualifications. More search cannot repair the 23 strict objective failures.
Ten of the 50 exact top-ten lists contain duplicate numerical/type mappings
realized by different mention indices. A separate 49-query follow-up ranks without evaluation values, conditional on the same supplied inventory. It reuses
40 already-distinct lists and performs 23 additional solves for the other
nine. Every retained top-ten list is certified, but criterion-correct list
retention remains 9/49, with zero gained or lost queries; selected correctness
remains 3/49. Thus reclaiming this alias budget supplies no observed
correctness gain and no justification for full-dataset expansion. Mention equality, numerical equality and semantic identity are
different concepts. No gold value is used to rank inference candidates.

Detailed source/reference conventions, the versioned correction overlay and the historical tolerance limitation are given in Appendix~\ref{app:reference-conventions}. No reference correction is used during inference, and no new hard unit or mathematical relation was validated by these studies.
